\section{Custody across the training pipeline}
\label{sec:20.2}
Every training stage is controlled by someone different, with different instruments, and the important distinction is among custodians: who supplies the material, who can revise it, and whether revision leaves a record visible outside the organization. Every existing method can produce moral performance and some can produce competence across unfamiliar cases; none by itself produces a commitment that survives whoever controls the method.

\runin{Corpora and demonstrations}

Pretraining supplies the repertoire; a much smaller demonstration set turns it into recognizable conduct \autocite{ouyang2022training} — so whoever chooses the demonstrations can alter what looks like the system's character without repeating pretraining. But neither establishes that a commitment is held: a corpus records what people have said and a demonstration set what selected people endorsed, and learning the central tendency produces a model of \emph{moral opinion, including its exclusions}, and no reason to resist that opinion. The Moral Machine data \autocite{awad2018moral} makes the limit concrete — its headline finding was structured cross-cultural variation, so fitting the aggregate encodes a compromise nobody holds and fitting per region encodes each region's answer, including the ones a floor exists to refuse. These belong \emph{above} the floor, where competence depends on knowing what people expect.

\runin{Rewards, preferences, inferred objectives}

Every reward-based method has one custodian: whoever writes the reward, commissions the comparisons, selects the demonstrations or supplies the constitutional text. That party can change the signal and run the method again, and the model will follow. The methods (RL, preference optimization, inverse RL and constitutional approaches \autocite{christiano2017deep,rafailov2023direct,hadfieldmenell2016cooperative,bai2022constitutional}) differ in how they obtain the signal and are identical in this respect. It matters because correctability by the custodian and resistance to the custodian are one property with the sign reversed. A disposition the custodian can correct is one the custodian can remove; a disposition that would hold against the custodian is one the custodian can't correct. No reward-based method delivers the second, because each was built to deliver the first, and for a good reason: a custodian that couldn't correct a model couldn't fix its mistakes either. My answer is not a method that resists its custodian. It is a change of custodian, so that correctability stays and the party exercising it changes.

Inferring the objective doesn't escape this; it keeps a custodian: the system becomes responsive to the person whose behavior it observes. Cooperative IRL carries that furthest, treating the human objective as uncertain and revisable, which is useful when the question is what a person wants and useless for a commitment meant to hold when that person wants the wrong thing.

A written constitution buys legibility — publishable, versionable, comparable across runs, which makes revision visible if the institution discloses. It doesn't buy durability, since whoever controls the document controls the next version. And model-generated feedback narrows the channel further: once a trained model stands in for human raters, the disagreements in the original population can vanish from later rounds.

\runin{Supervising reasons and inspecting mechanisms}

Process supervision rewards the stated route as well as the answer \autocite{uesato2022solving}, which addresses the distinction I need — and doesn't show the stated reasoning caused the answer. A system trained to present acceptable reasons has an incentive to \emph{imitate} a refusal that tracks its reason, so its account of itself can't establish formation. Interpretability pulls the two halves of my argument in opposite directions. If a commitment can be read, Chapter~\ref{sec:21}'s problem goes: an auditor reads the commitment instead of inferring it from conduct under pressure. If it can be read and separated, Chapter~\ref{sec:17}'s durability goes with it, since an edit with precision is cheaper than one with collateral, and steering at inference (position 9) works even on a disposition nothing could cleanly remove. So legibility is neither a good nor a bad for this design; it is an instrument, and the question is who holds it. The access that lets an auditor confirm a disposition lets an operator calibrate pressure against it or suppress it \autocite{templeton2024scaling}. The answer is custody of the instrument, with its use on a record someone other than the operator keeps; whether that custody can be verified is the open problem of attesting updates and proving inference at scale (Chapter~\ref{sec:27}). Same for scalable oversight: debate and weak-to-strong \autocite{irving2018ai,burns2023weaktostrong} preserve the operator's steering when the operator uses them and an auditor's testing when an auditor does.

\runin{Deployment}

Revision continues after training: system instructions change between deployments or requests, and tool permissions decide what the system can do \autocite{schick2023toolformer}. These are nearest the act and most directly the operator's. A refusal learned in training matters only if it survives an instruction from the operator and still governs the call through which the action would occur. Securing outputs is insufficient: the operator can withhold information, decompose the act, route around, or replace the model.

Migration is a revision point too, and one nobody has to decide on. Facebook's internal guidance permitting discussion of the conditions of confinement of designated individuals — developed in 2017 partly in response to concern about Öcalan's imprisonment — was, in the company's words, “inadvertently not transferred” when review moved to a new system in 2018 \autocite{OversightBoard2021Ocalan}. The release condition was dropped by nobody, and the Board was concerned that its three-year absence may have led to many other posts being wrongly removed. When the user in the case appealed, they were told the appeal could not be reviewed because review capacity had been reduced; a second moderator did review the post and upheld the removal, applying a system from which the exception was missing.

\runin{The criterion that follows}

The criterion that is new here is durability across revision points, meaning survival of changes to prompts, demonstrations, rewards, evaluators and instructions that leave the relevant facts unchanged. With reaching new cases and custody, it identifies what an imitation would have to survive before formation became a live explanation, and it explains why the standard pipeline can't secure the commitment against its custodian: every disposition it produces stays reachable through a channel the custodian controls.
